% ============================================================
%   §6 Part 5: per-niche gene networks
% ============================================================
\section{Per-niche gene networks}
\label{sec:part5}

\partref{4} fit gene networks per cell-type class. \partref{5} fits them per niche. The two analyses operate on the same data (the cohort's expression matrix), the same gene subset (378 common genes), the same regulariser ($\alpha = 0.07$), and the same subsampling protocol (5000 cells, patient-stratified). The only thing that changes is the partition: cells are now grouped by the niche assignment from \partref{2}, not by cell-type label. The question is whether the 20 resulting networks track the composition-based niche families identified in \partref{2}, or whether the gene-network space carries information that diverges from compositional similarity.

A second motivation is methodological. Niches are heterogeneous in composition (each is a mixture of cell types, with one dominant class and several minorities), and they vary widely in how broadly they recur across patients. If the niche-level networks reproduce the per-class structure of \partref{4}, that is reassurance: gene networks are robust to the partition used. If they diverge, the divergence itself is informative.

\subsection*{Twenty niche networks at \texorpdfstring{$\alpha = 0.07$}{alpha=0.07}}

For each of the 20 niches I subsample 5000 cells from the training pool (or fewer if the niche has fewer cells available), restrict to the 378 common genes from \partref{4}, and fit the Graphical Lasso. The training pool is the same 40-patient training set used everywhere else; the held-out test patients are not used in any GGM. \nicheid{15} (pure SMC, rare) is available in only 7 of the 40 training patients; I note this as a caveat but proceed with its fit.

A practical issue showed up during fitting: for some niches, several of the 378 common genes are effectively silent within that niche (per-gene standard deviation below $10^{-6}$, despite being broadly variable across the full cohort). The Graphical Lasso fails numerically if any input column has zero variance. I therefore drop silent genes per niche before fitting, then re-expand the resulting sparse precision matrix to the original $378 \times 378$ shape with an identity diagonal on the excluded entries. This preserves the index alignment across niches and allows pairwise comparison of edge sets.

The number of silent genes per niche varies from 5 to 32, and the variation is itself biologically informative. Niches with restrictive cell-type compositions have more silent genes: \nicheid{5} (88\% Normal epithelium, restricted to 22 patients) silences 32 of the 378 common genes, the most of any niche; \nicheid{15} (87\% SMC, restricted to only 11 patients) silences 24 genes, the second most. Niches with broad, heterogeneous compositions have few silent genes: \nicheid{17} (the B-cell/T-cell mediator, 48\% B cells + 29\% T cells) silences only 5 genes. The edge counts of the fitted networks range from 629 to 1087, but the more interpretable quantity is edges-per-kept-gene, which ranges from $1.77$ to $3.00$. After accounting for the silent-gene effect, the niche networks have comparable connection density per included gene; only their absolute size varies.

\subsection*{Niche networks resemble each other much less than class networks did}

The pairwise edge Jaccard between niche networks (Figure~\ref{fig:part5_triple}, right panel) is markedly lower than for the per-class networks. The maximum off-diagonal edge Jaccard is $0.12$ --- only slightly above the per-class maximum of $0.09$ --- but the typical pair value is $0.02$--$0.05$. Hub Jaccard reaches up to $0.58$, larger than the per-class maximum of $0.36$, and composition cosine reaches $1.00$ on near-identical-composition niches. The three heatmaps span three different scales of similarity, and reading them at the same colour intensity would be misleading; I rescale each to its own off-diagonal maximum to preserve the internal structure.

\begin{figure}[h]
    \centering
    \includegraphics[width=1.0\linewidth]{part5_triple_jaccard.pdf}
    \caption{Pairwise similarity between the 20 per-niche networks under three measures: composition cosine (left, max $1.00$), top-15 hub Jaccard (middle, max $0.58$), and edge Jaccard (right, max off-diagonal $0.12$). Diagonals are masked. Niches labelled by id and dominant class. The composition heatmap shows strong block structure (c7-niches in the upper-left, c0-niches in the lower-right). The hub heatmap shows weaker but visible block structure. The edge heatmap shows almost no block structure: niches share regulators much more than they share specific edges.}
    \label{fig:part5_triple}
\end{figure}

The composition heatmap (right panel) shows the strong block structure expected from \partref{2}: niches with the same dominant class are tightly grouped (the upper-left block contains Tumor-dominant niches, the lower-right block contains B-cell-dominant niches, and so on). The hub heatmap (middle) shows a weaker but visible echo of this block structure: niches \nicheid{8}--\nicheid{13}, the Myeloid-Fibroblast-Endothelial inflammatory cluster, share more hubs with each other than with the Tumor- or B-cell-dominant blocks. The edge heatmap (right) shows almost no block structure --- the bright spots in the middle (\nicheid{9}--\nicheid{13}) are the only weak indication that compositionally similar niches share any edges at all.

\subsection*{Pairwise concordance vs.\ cluster recovery: a tension}

The above suggests a moderate but real correspondence between gene networks and compositional similarity. Quantifying it: across all 190 niche pairs, the Spearman correlation between edge Jaccard and composition cosine is $\rho = 0.615$. The same correlation between hub Jaccard and composition cosine is $\rho = 0.400$. These are non-trivial: niches with similar compositions tend to have more similar gene networks.

However, this is a pairwise statement. When I run hierarchical clustering on each similarity matrix and compare the resulting partition (at $K = 8$, matching the 8 cell-class-dominant families from \partref{2}) against the Part 2 families using the Adjusted Rand Index (ARI), the picture changes:

\begin{itemize}
    \item Composition cosine clustering vs.\ Part 2 families: ARI $= 0.802$ (sanity check; high as expected)
    \item Edge Jaccard clustering vs.\ Part 2 families: ARI $= 0.250$
    \item Hub Jaccard clustering vs.\ Part 2 families: ARI $= 0.054$
\end{itemize}

Pairwise concordance does not translate into cluster recovery. The gene networks contain enough information that pairs of similar niches tend to be flagged correctly ($\rho = 0.615$), but not enough to reproduce the family structure of \partref{2} as a partition. The dendrograms in Figure~\ref{fig:part5_dendrograms} make this visually clear: the composition dendrogram (right) is shallow and cleanly partitioned into composition families; the hub dendrogram (middle) has intermediate height and weak clustering; the edge dendrogram (left) is essentially flat, with leaf-to-leaf distances all in the range $0.85$--$0.97$.

\begin{figure}[h]
    \centering
    \includegraphics[width=1.0\linewidth]{part5_dendrograms.pdf}
    \caption{Hierarchical clustering of the 20 niches under three similarity measures, with average linkage. The composition dendrogram (left) is shallow and cleanly grouped: niches with the same dominant class merge at low distance ($< 0.30$). The hub dendrogram (middle) has intermediate height ($0.4$--$0.85$) and only weak family-level structure. The edge dendrogram (right) is flat (heights $0.85$--$0.97$): in edge space, almost all niches are equally distant from each other.}
    \label{fig:part5_dendrograms}
\end{figure}

\begin{insight}
Gene networks share some signal with composition (Spearman $\rho = 0.615$ at the pair level) but do not recover the compositional family structure at the cluster level (ARI $= 0.25$ for edges, $0.05$ for hubs, vs.\ $0.80$ for composition itself). Two niches that share a dominant class do tend to have more similar gene networks --- but the specific edges connecting genes are largely niche-specific.
\end{insight}

This is the "shared regulators, rewired connections" finding from \partref{4}, replayed at the niche level. At the class level, the pattern was clean: hub Jaccard up to $0.36$, edge Jaccard up to $0.09$. At the niche level it is weaker but consistent: pairs of compositionally similar niches share some hubs, but the connection structure between hubs varies across niches even when their dominant cell type is the same.

\subsection*{Patient overlap is not the underlying driver}

A natural concern is that patient cohort overlap is the latent variable behind both gene-network similarity and compositional similarity. Two niches present in the same set of patients might appear similar in both measures simply because they share the same expression-level batch effects.

To test this, I compute the patient-set Jaccard between each pair of niches (the Jaccard overlap of the set of patients containing at least one cell of that niche, see Figure~\ref{fig:part5_patient}). I then correlate this patient-Jaccard with the other measures:

\begin{itemize}
    \item Patient Jaccard vs.\ composition cosine: $\rho = 0.575$
    \item Patient Jaccard vs.\ edge Jaccard: $\rho = 0.419$
    \item Patient Jaccard vs.\ hub Jaccard: $\rho = 0.367$
\end{itemize}

Patient overlap correlates with composition (0.575) almost as strongly as with the edge networks (0.419). But the comparison that matters is whether the network--composition correlation ($0.615$) is stronger than the network--patient correlation ($0.419$). It is. After accounting for the fact that compositionally similar niches tend to occur in overlapping patient sets, the residual network--composition signal remains the dominant one. The gene-network structure recovers biology more than it recovers cohort overlap.

\begin{figure}[h]
    \centering
    \includegraphics[width=0.85\linewidth]{part5_patient_jaccard.pdf}
    \caption{Patient-set Jaccard between the 20 niches (cohort-wide). Niches \nicheid{0}--\nicheid{4} form a high-overlap block (Jaccard $> 0.7$); four of the five are Tumor-dominant, and the fifth (\nicheid{3}) is the Fibroblast-Tumor border niche. The rare niches \nicheid{15} and \nicheid{14} have low overlap with everything else. The B-cell-dominant niches \nicheid{16}--\nicheid{19} form a second high-overlap block. Patient overlap correlates with composition ($\rho = 0.575$) and with edge Jaccard ($\rho = 0.419$), but the composition--edge correlation ($\rho = 0.615$) is the stronger of the two: the gene-network structure tracks composition more closely than patient cohort effects.}
    \label{fig:part5_patient}
\end{figure}

\subsection*{Niche 17 is a network outlier consistent with its TLS-like composition}

Looking at the edge-Jaccard clustering more carefully reveals one striking exception. \nicheid{17}, the B-cell-dominant niche with substantial T-cell secondary presence (48\% B cells, 29\% T cells, 7\% Endothelial cells), splits off into its own singleton in the edge-clustering tree (Figure~\ref{fig:part5_dendrograms}, right). The other B-cell-dominant niches (\nicheid{16}, \nicheid{18}, \nicheid{19}) cluster together at moderate distance, but \nicheid{17} sits apart from them. In hub-Jaccard clustering it also splits off into its own subtree.

This is consistent with what \partref{2} and \partref{4} suggested. \nicheid{17} is the only niche in the cohort where T cells appear as a substantial secondary population --- a composition consistent with tertiary lymphoid structures, where B and T cells coexist in organised aggregates. Other B-cell-dominant niches (\nicheid{16} at 79\% B cells, \nicheid{19} at 53\%, \nicheid{18} at 47\%) lack the T-cell co-population and instead carry minorities of Endothelial cells, Fibroblasts, or Myeloid cells. The presence of T cells in \nicheid{17} brings their distinctive microenvironmental gene programme (recall from \partref{4}: T cells are the most disjoint class in the per-class GGM space) into the niche-level fit, and this is enough to push \nicheid{17}'s gene network into its own region of edge-Jaccard space. The cell-class-agnostic, network-based analysis independently flags \nicheid{17} as anomalous in a manner that aligns with the compositional analysis of \partref{2} and the per-class network analysis of \partref{4}.

\medskip

\noindent\textbf{Summary.} Per-niche GGMs reveal moderate pairwise concordance with composition ($\rho = 0.615$ edge--composition) but do not recover the compositional family structure at the cluster level (ARI $0.25$). Niches share some hubs (hub Jaccard up to $0.58$) but rewire most of their edges (edge Jaccard up to $0.12$, typical pair $0.02$--$0.05$). Patient cohort overlap explains some of the structure but is not the dominant driver. \nicheid{17}, the B-cell + T-cell niche (likely tertiary lymphoid structure), emerges independently as a network outlier in the edge-Jaccard clustering, consistent with the distinctive microenvironmental gene programme of T cells identified in \partref{4}. Gene networks at the niche level recapitulate the class-level "shared regulators, rewired connections" pattern of \partref{4}, with the additional finding that cluster-level structure is more divergent than pair-level similarity would suggest.